% Part: normal-modal-logic
% Chapter: tableaux
% Section: countermodels

\documentclass[../../../include/open-logic-section]{subfiles}

\begin{document}

\olfileid[pt]{nml}{tab}{cou}
      
\olsection{Contramodelos a partir de \usetoken{P}{tableau}}

\begin{explain}
  A demonstração do teorema da completude não mostra apenas que, se
  $\Entails !A$, então $\Proves !A$; ela também nos fornece um método
  para construir contramodelos para~$!A$ se $\Entails/ A$. No caso
  de~$\Log{K}$, esse método constitui um \emph{procedimento de decisão}.
  Com efeito, suponha $\Entails/ !A$. A demonstração de
  \olref[cpl]{prop:complete-tableau} fornece então um método para construir
  um !!{tableau} completo. Na verdade, o método sempre termina. As regras
  proposicionais para~$\Log{K}$ acrescentam apenas !!{formula}s prefixadas
  de menor complexidade, O.s.\@, cada regra proposicional só precisa ser
  aplicada uma vez em um ramo para cada fórmula assinalada
  $\sFmla{S}{!A}[\sigma]$. Prefixos novos são gerados apenas pela
  \iftag{prvBox}{$\TRule{\False}{\Box}$}{}\iftag{notprvBox,notprvDiamond}{}{
    e }\iftag{prvDiamond}{$\TRule{\True}{\Diamond}$}{}
  \iftag{notprvBox,notprvDiamond}{regra}{regras}, que também só precisa
  ser aplicada uma vez (e produz um único prefixo novo).
  \iftag{prvBox}{$\TRule{\True}{\Box}$}{}\iftag{notprvBox,notprvDiamond}{}{
    e }\iftag{prvDiamond}{$\TRule{\False}{\Diamond}$}{}
  \iftag{notprvBox,notprvDiamond}{pode}{podem} precisar ser aplicada
  várias vezes, mas apenas uma vez por prefixo, e somente um número finito
  de prefixos novos é gerado. Assim, depois de um número finito de etapas,
  a construção produz um ramo fechado ou um ramo completo.

  Uma vez construído um tableau com um ramo aberto completo, a demonstração
  de \olref[cpl]{thm:tableau-completeness} nos fornece um modelo explícito
  que satisfaz o conjunto original de !!{formula}s prefixadas. Portanto,
  não apenas ocorre que, se $\Gamma \Entails !A$, então existe um
  !!{tableau} fechado e $\Gamma \Proves !A$: se procurarmos o !!{tableau}
  fechado da maneira correta e terminarmos com um !!{tableau} ``completo'',
  não só saberemos que $\Gamma \Entails/ !A$, mas poderemos de fato
  construir um contramodelo.
\end{explain}

\iftag{prvBox}{
\begin{ex}
  Sabemos que $\Proves/ \Box(p \lor q) \lif (\Box p \lor \Box q)$.
  A construção de um tableau começa com:
  \begin{oltableau}
    [\pFmla{\False}{\Box(p \lor q) \lif (\Box p \lor \Box q)}{1},
      just = \TAss, checked
      [\pFmla{\True}{\Box(p \lor q)}{1},
        just = {\TRule{\False}{\lif}[1]}, 
        [\pFmla{\False}{\Box p \lor \Box q}{1},
          just = {\TRule{\False}{\lif}[1]}, checked
          [\pFmla{\False}{\Box p}{1},
            just = {\TRule{\False}{\lor}[3]}, checked
            [\pFmla{\False}{\Box q}{1},
              just = {\TRule{\False}{\lor}[3]}, checked
              [\pFmla{\False}{p}{1.1}, 
                just = {\TRule{\False}{\Box}[4]}, checked
                [\pFmla{\False}{q}{1.2}, 
                  just = {\TRule{\False}{\Box}[5]}, checked
                ]
              ]
            ]
          ]
        ]
      ]
    ]
  \end{oltableau}
  É claro que o !!{tableau} ainda não está terminado. Na etapa seguinte,
  consideramos a única linha sem marca de verificação: a !!{formula}
  prefixada $\sFmla{\True}{\Box(p \lor q)}[1]$ na linha~$2$. A construção
  do tableau fechado determina a aplicação da regra
  $\TRule{\True}{\Box}$ a todo prefixo usado no ramo, O.s.\@, tanto a
  $1.1$ quanto a $1.2$:
  \begin{oltableau}
    [\pFmla{\False}{\Box(p \lor q) \lif (\Box p \lor \Box q)}{1},
      just = \TAss, checked
      [\pFmla{\True}{\Box(p \lor q)}{1},
        just = {\TRule{\False}{\lif}[1]}, 
        [\pFmla{\False}{\Box p \lor \Box q}{1},
          just = {\TRule{\False}{\lif}[1]}, checked
          [\pFmla{\False}{\Box p}{1},
            just = {\TRule{\False}{\lor}[3]}, checked
            [\pFmla{\False}{\Box q}{1},
              just = {\TRule{\False}{\lor}[3]}, checked
              [\pFmla{\False}{p}{1.1}, 
                just = {\TRule{\False}{\Box}[4]}, checked
                [\pFmla{\False}{q}{1.2}, 
                  just = {\TRule{\False}{\Box}[5]}, checked
                  [\pFmla{\True}{p \lor q}{1.1}, 
                    just = {\TRule{\True}{\Box}[2]}
                    [\pFmla{\True}{p \lor q}{1.2}, 
                      just = {\TRule{\True}{\Box}[2]}
                    ]
                  ]
                ]
              ]
            ]
          ]
        ]
      ]
    ]
  \end{oltableau}
  Agora as linhas 2, 8 e 9 não têm marcas de verificação. Mas nenhum
  prefixo novo foi acrescentado; por isso, aplicamos
  $\TRule{\True}{\lor}$ às linhas~8 e~9, em todos os ramos resultantes
  (desde que não fechem):
  \begin{oltableau}
    [\pFmla{\False}{\Box(p \lor q) \lif (\Box p \lor \Box q)}{1},
      just = \TAss, checked
      [\pFmla{\True}{\Box(p \lor q)}{1},
        just = {\TRule{\False}{\lif}[1]}, checked
        [\pFmla{\False}{\Box p \lor \Box q}{1},
          just = {\TRule{\False}{\lif}[1]}, checked
          [\pFmla{\False}{\Box p}{1},
            just = {\TRule{\False}{\lor}[3]}, checked
            [\pFmla{\False}{\Box q}{1},
              just = {\TRule{\False}{\lor}[3]}, checked
              [\pFmla{\False}{p}{1.1}, 
                just = {\TRule{\False}{\Box}[4]}, checked
                [\pFmla{\False}{q}{1.2}, 
                  just = {\TRule{\False}{\Box}[5]}, checked
                  [\pFmla{\True}{p \lor q}{1.1}, 
                    just = {\TRule{\True}{\Box}[2]}, checked
                    [\pFmla{\True}{p \lor q}{1.2}, 
                      just = {\TRule{\True}{\Box}[2]}, checked
                      [\pFmla{\True}{p}{1.1},
                        just = {\TRule{\True}{\lor}[8]}, checked, close
                      ]
                      [\pFmla{\True}{q}{1.1},
                        just = {\TRule{\True}{\lor}[8]}, checked
                        [\pFmla{\True}{p}{1.2},
                          just = {\TRule{\True}{\lor}[9]}, checked]
                        [\pFmla{\True}{q}{1.2},
                          just = {\TRule{\True}{\lor}[9]}, checked, close]
                      ]
                    ]
                  ]
                ]
              ]
            ]
          ]
        ]
      ]
    ]
  \end{oltableau}
  Resta um único ramo aberto, e ele é completo. A partir dele, definimos
  o modelo com os mundos $W = \{1, 1.1, 1.2\}$ (os únicos prefixos que
  aparecem no ramo aberto), a relação de acessibilidade
  $R = \{\tuple{1, 1.1}, \tuple{1, 1.2}\}$ e a atribuição
  $V(p) = \{1.2\}$ (pois a linha~11 contém
  $\sFmla{\True}{p}[1.2]$) e $V(q) = \{1.1\}$ (pois a linha~10 contém
  $\sFmla{\True}{q}[1.1]$). O modelo é representado em
  \olref{fig:counter-Box}, e pode-se verificar que ele é um contramodelo
  para $\Box(p \lor q) \lif (\Box p \lor \Box q)$.
  \begin{figure}
  \begin{center}
    \begin{tikzpicture}[modal]
      \node[world] (w1) [label={[align=right]right:\mFalse{p}\\ \mFalse{q}}]{$1$}; 
      \node[world] (w2) [label={[align=right]right:\mFalse{p}\\ \mTrue{q}},
        above left=of w1]{$1.1$}; 
      \node[world] (w3) [label={[align=right]right:\mTrue{p}\\ \mFalse{q}},
        above right=of w1] {$1.2$};
      \draw[->] (w1) to (w2);
      \draw[->] (w1) to (w3);
    \end{tikzpicture}
  \end{center}
  \caption{Um contramodelo para $\Box(p \lor q) \lif (\Box p \lor \Box
    q)$.}
  \ollabel{fig:counter-Box}
\end{figure}
\end{ex}          
}
{
\begin{ex}
  Sabemos que $\Proves/ (\Diamond p \land \Diamond q) \lif \Diamond(p
  \land q)$. A construção de um tableau começa por:
  \begin{oltableau}
    [\pFmla{\False}{(\Diamond p \land \Diamond q) \lif \Diamond(p \land q)}{1},
      just = \TAss, checked
      [\pFmla{\True}{\Diamond p \land \Diamond q}{1},
        just = {\TRule{\False}{\lif}[1]}, checked
        [\pFmla{\False}{\Diamond(p \land q)}{1},
          just = {\TRule{\False}{\lif}[1]}
          [\pFmla{\True}{\Diamond p}{1},
            just = {\TRule{\True}{\land}[2]}, checked
            [\pFmla{\True}{\Diamond q}{1},
              just = {\TRule{\True}{\land}[2]}, checked
              [\pFmla{\True}{p}{1.1}, 
                just = {\TRule{\True}{\Diamond}[4]}, checked
                [\pFmla{\True}{q}{1.2}, 
                  just = {\TRule{\True}{\Diamond}[5]}, checked
                ]
              ]
            ]
          ]
        ]
      ]
    ]
  \end{oltableau}
  É claro que o !!{tableau} ainda não está terminado. Na etapa seguinte,
  consideramos a única linha sem marca de verificação: a !!{formula}
  prefixada $\sFmla{\True}{\Diamond(p \land q)}[1]$ na linha~$3$. A construção
  do tableau fechado determina a aplicação da regra
  $\TRule{\True}{\Diamond}$ a todo prefixo usado no ramo, O.s.\@, tanto a
  $1.1$ quanto a $1.2$:
  \begin{oltableau}
    [\pFmla{\False}{\Diamond(p \land q) \lif (\Diamond p \land \Diamond q)}{1},
      just = \TAss, checked
      [\pFmla{\True}{\Diamond p \land \Diamond q}{1},
        just = {\TRule{\False}{\lif}[1]}, checked
        [\pFmla{\False}{\Diamond (p \land q)}{1},
          just = {\TRule{\False}{\lif}[1]}
          [\pFmla{\True}{\Diamond p}{1},
            just = {\TRule{\True}{\land}[2]}, checked
            [\pFmla{\True}{\Diamond q}{1},
              just = {\TRule{\True}{\land}[2]}, checked
              [\pFmla{\True}{p}{1.1}, 
                just = {\TRule{\True}{\Diamond}[4]}, checked
                [\pFmla{\True}{q}{1.2}, 
                  just = {\TRule{\True}{\Diamond}[5]}, checked
                  [\pFmla{\False}{p \land q}{1.1}, 
                    just = {\TRule{\False}{\Diamond}[3]}
                    [\pFmla{\False}{p \land q}{1.2}, 
                      just = {\TRule{\False}{\Diamond}[3]}
                    ]
                  ]
                ]
              ]
            ]
          ]
        ]
      ]
    ]
  \end{oltableau}
  Agora as linhas 3, 8 e 9 não têm marcas de verificação. Mas nenhum
  prefixo novo foi acrescentado; por isso, aplicamos
  $\TRule{\False}{\land}$ às linhas~8 e~9, em todos os ramos resultantes
  (desde que não fechem):
  \begin{oltableau}
    [\pFmla{\False}{(\Diamond p \land \Diamond q) \lif \Diamond(p \land q)}{1},
      just = \TAss, checked
      [\pFmla{\True}{\Diamond p \land \Diamond q}{1},
        just = {\TRule{\False}{\lif}[1]}, checked
        [\pFmla{\False}{\Diamond(p \land q)}{1},
          just = {\TRule{\False}{\lif}[1]}
          [\pFmla{\True}{\Diamond p}{1},
            just = {\TRule{\True}{\land}[2]}, checked
            [\pFmla{\True}{\Diamond q}{1},
              just = {\TRule{\True}{\land}[2]}, checked
              [\pFmla{\True}{p}{1.1}, 
                just = {\TRule{\True}{\Diamond}[4]}, checked
                [\pFmla{\True}{q}{1.2}, 
                  just = {\TRule{\True}{\Diamond}[5]}, checked
                  [\pFmla{\False}{p \land q}{1.1}, 
                    just = {\TRule{\False}{\Diamond}[3]}, checked
                    [\pFmla{\False}{p \land q}{1.2}, 
                      just = {\TRule{\False}{\Diamond}[3]}, checked
                      [\pFmla{\False}{p}{1.1},
                        just = {\TRule{\False}{\land}[8]}, checked, close
                      ]
                      [\pFmla{\False}{q}{1.1},
                        just = {\TRule{\False}{\land}[8]}, checked
                        [\pFmla{\False}{p}{1.2},
                          just = {\TRule{\False}{\land}[9]}, checked]
                        [\pFmla{\False}{q}{1.2},
                          just = {\TRule{\False}{\land}[9]}, checked, close]
                      ]
                    ]
                  ]
                ]
              ]
            ]
          ]
        ]
      ]
    ]
  \end{oltableau}
  Resta um único ramo aberto, e ele é completo. A partir dele, definimos
  o modelo com os mundos $W = \{1, 1.1, 1.2\}$ (os únicos prefixos que
  aparecem no ramo aberto), a relação de acessibilidade
  $R = \{\tuple{1, 1.1}, \tuple{1, 1.2}\}$ e a atribuição $V(p) =
  \{1.1\}$ (pois a linha~6 contém $\sFmla{\True}{p}[1.1]$) e
  $V(q) = \{1.2\}$ (pois a linha~7 contém
  $\sFmla{\True}{q}[1.1]$). O modelo é representado em
  \olref{fig:counter-Diamond}, e pode-se verificar que ele é um
  contramodelo para $(\Diamond p \land \Diamond q) \lif \Diamond (p \land q)$.
  \begin{figure}
  \begin{center}
    \begin{tikzpicture}[modal]
      \node[world] (w1) [label={[align=right]right:\mFalse{p}\\ \mFalse{q}}]{$1$}; 
      \node[world] (w2) [label={[align=right]right:\mTrue{p}\\ \mFalse{q}},
        above left=of w1]{$1.1$}; 
      \node[world] (w3) [label={[align=right]right:\mFalse{p}\\ \mTrue{q}},
        above right=of w1] {$1.2$};
      \draw[->] (w1) to (w2);
      \draw[->] (w1) to (w3);
    \end{tikzpicture}
  \end{center}
  \caption{Um contramodelo para $(\Diamond p \land \Diamond q) \lif
    \Diamond (p \land q)$.}
  \ollabel{fig:counter-Diamond}
\end{figure}
\end{ex}          
}

\end{document}
